\clearpage
\section{The Pomodoro cycle, end to end}

\subsection{Normal case: a 4-session cycle}

\begin{figure}[H]
\centering
\begin{tikzpicture}[x=1mm,y=1mm]
  % nine blocks: S B S B S B S B LB
  \foreach \i/\lab/\col/\txt in {
    0/S1/pastel!45/{25m},
    1/B1/mint!55/{5m},
    2/S2/pastel!45/{25m},
    3/B2/mint!55/{5m},
    4/S3/pastel!45/{25m},
    5/B3/mint!55/{5m},
    6/S4/pastel!45/{25m},
    7/LB/mandarin!45/{15m},
    8/S1/pastel!45/{25m}}
  {
    \pgfmathsetmacro{\x}{\i*16}
    \fill[\col] (\x,30) rectangle ++(13,16);
    \node[anchor=center,font=\tiny\sffamily\bfseries,text=ink] at ({\x+6.5},40) {\txt};
    \node[anchor=center,font=\scriptsize\sffamily,text=softink] at ({\x+6.5},34) {\lab};
    \draw[e] ({\x+13},38) -- ({\x+15.4},38);
  }
  % the cycle wraps
  \draw[ethin, softink!50] (8*16+6.5,30) -- (8*16+6.5,24) -- (6.5,24) -- (6.5,29.4);
  \node[anchor=west,font=\tiny\sffamily,text=softink] at (0,18)
    {\texttt{completed\_sessions} = 4 \quad and \quad 4 mod 4 = 0 \quad$\Rightarrow$\quad \textbf{LongBreak}};
  % pips
  \foreach \i in {0,1,2,3} {
    \fill[mint!60!black] ({2+\i*32},8) circle (2.2);
  }
  \node[anchor=west,font=\tiny\sffamily,text=softink] at (0,1)
    {pips: \texttt{completed\_sessions \% long\_break\_cycle}};
\end{tikzpicture}
\caption{With the default settings (\texttt{long\_break\_cycle = 4}) the fourth finished
session earns a long break. Skipping any session breaks the chain --- the counter is
only incremented by \texttt{Completion::Finished}.}
\label{fig:cycle}
\end{figure}

\subsection{The three switches that reshape the cycle}

\begin{figure}[H]
\centering
\begin{tikzpicture}[
  sw/.style={rounded corners=2pt,draw=ink!50,fill=white,align=center,inner sep=5pt,
             font=\scriptsize\sffamily,text width=44mm,minimum height=12mm},
]
  \node[sw, draw=accent!60, fill=accent!10, text=accent!55!black] (d1) at (0,0)
    {\textbf{disable\_breaks = true}\\[1.5mm]{\scriptsize after \emph{every} session the next\\activity is
     \texttt{Session}. \texttt{completed\_sessions} still increments.}};

  \node[sw, draw=mandarin!70!black, fill=mandarin!16] (d2) at (0,-3.0)
    {\textbf{disable\_long\_breaks = true}\\[1.5mm]{\scriptsize the cycle counter keeps running, but
     \texttt{due\_long} is ignored, so every break is a short break.}};

  \node[sw, draw=mint!70!black, fill=mint!18] (d3) at (0,-6.0)
    {\textbf{long\_break\_cycle = $n$}\\[1.5mm]{\scriptsize the long break moves to session $n$, $2n$,
     $3n$, \dots \texttt{.max(1)} guards a zero read from a hand-edited file.}};

  \node[sw, draw=pastel!45!black, fill=corefill] (d4) at (6.6,-3.0)
    {\textbf{show\_completed\_only}\\[1.5mm]{\scriptsize not a timer switch --- it filters the
     \emph{statistics} view so skipped activities are excluded from the totals.}};
\end{tikzpicture}
\caption{Four settings that change the shape of the cycle. The first three are read
inside \texttt{advance()} on every transition; the fourth only affects the numbers shown
in the statistics panel.}
\label{fig:switches}
\end{figure}

\subsection{Auto-start, and who decides it}

\begin{lstlisting}[caption={\file{app.rs}, \texttt{complete\_current()}. Note that the
core decides \emph{whether} to start the next activity, but never starts it itself
unless asked.}]
pub fn complete_current(&mut self) {
    let finished = self.timer.activity();
    self.status = format!("{} finished", finished.label());
    self.record(true);                       // write the CSV row first
    self.timer.finish();                     // advance; state becomes Idle
    self.notified = false;
    let auto = if self.timer.activity().is_break() {   // note: the NEXT activity
        self.settings.auto_start_break
    } else {
        self.settings.auto_start_session
    };
    if auto { self.timer.start(); }
}
\end{lstlisting}

\begin{figure}[H]
\centering
\begin{tikzpicture}[
  b/.style={rounded corners=2pt,draw=ink!50,fill=white,align=center,inner sep=5pt,
            font=\scriptsize\sffamily,text width=40mm,minimum height=11mm},
  k/.style={b,draw=accent!65,fill=accent!12,text=accent!55!black},
  g/.style={b,draw=mint!70!black,fill=mint!18,text=mint!60!black},
]
  \node[b] (t) {the displayed second\\reaches zero};
  \node[k] (n) {\texttt{tick()} raises\\\texttt{notify\_requested} \emph{once}};
  \node[b] (o) {\textbf{unless} \texttt{overtime},\\\texttt{complete\_current()}};
  \node[b] (s) {look up\\\texttt{auto\_start\_break}\\\emph{or}\\\texttt{auto\_start\_session}};
  \node[g] (st) {\texttt{timer.start()}\\[-0.5mm]{\scriptsize the next activity begins immediately}};

  \draw[e] (t)--(n); \draw[e] (n)--(o); \draw[e] (o)--(s); \draw[e] (s)--(st);
  \draw[e] (o.east) -- ++(7mm,0);
  \node[note, anchor=west, text width=54mm, align=left, font=\tiny\sffamily] at (7.4,-0.55)
    {\emph{With overtime on, the chain stops here.}\\
     \texttt{timer} keeps running, the display shows \texttt{+MM:SS},\\
     and the user ends it with \texttt{Finish} or the \texttt{F} key.};
\end{tikzpicture}
\caption{End-of-activity path with overtime off. With \texttt{overtime} on, the
completion is deferred and the activity stays running past zero.}
\label{fig:finishpath}
\end{figure}

\subsection{Overtime: the display flips sign}

\begin{lstlisting}[caption={\file{app.rs}. Overtime is a \emph{display} mode plus a
deferred completion; the timer itself never changes state.}]
pub fn display_seconds(&self) -> u64 {
    if self.is_overtime() {
        self.timer.elapsed().saturating_sub(self.timer.duration()).as_secs()
    } else {
        self.timer.remaining_display().as_secs()
    }
}
pub fn is_overtime(&self) -> bool {
    self.settings.overtime
        && self.timer.remaining().is_zero()
        && self.timer.state() != RunState::Idle
}
\end{lstlisting}

\begin{figure}[H]
\centering
\begin{tikzpicture}[x=1mm,y=1mm]
  \draw[draw=softink!30,line width=0.5pt] (0,0) rectangle (150,34);
  \draw[-stealth, draw=ink!60] (6,17) -- (144,17);
  \fill[pastel!50] (6,12) rectangle (86,15);
  \fill[mint!50] (86,12) rectangle (96,15);
  \fill[accent!45] (96,12) rectangle (142,15);
  \node[anchor=north,font=\tiny\sffamily,text=softink] at (6,10) {countdown};
  \node[anchor=north,font=\tiny\sffamily,text=mint!60!black] at (91,10) {zero};
  \node[anchor=north,font=\tiny\sffamily,text=accent!55!black] at (119,10) {overtime: display reads \texttt{+MM:SS}};
  \node[anchor=south,font=\tiny\sffamily,text=softink] at (6,28)
    {\texttt{tick()} stops calling \texttt{complete\_current()} while overtime is on};
  \draw[er] (91,20) -- (91,27);
  \node[anchor=west,align=left,font=\tiny\sffamily,text=accent!55!black] at (94,25)
    {a \texttt{Finish} button appears\\in both shells --- nothing else changes};
  \node[anchor=south,font=\tiny\sffamily,text=softink] at (6,3)
    {\texttt{remaining()} is \texttt{duration.saturating\_sub(elapsed())}: it pins at zero, so the
     transition is smooth};
\end{tikzpicture}
\caption{Overtime is a display convention, not a fourth \texttt{RunState}.}
\label{fig:overtime}
\end{figure}